\section{Metodologia}

\begin{frame}{Visão Geral da Metodologia}
  \textbf{Caracterização:} Relato de experiência (Estudo de caso prático no SFM).

  \vspace{1em}

  \begin{figure}
      \centering
      \caption{Fases de execução da metodologia (não linear)}

      \vspace{-1em}
    \begin{tikzpicture}[
      node distance=1.5cm and 2cm,
      stepbox/.style={rectangle, rounded corners, draw=black, thick, fill=gray!10, minimum width=4.5cm, minimum height=1.5cm, text centered, text width=3.5cm},
      arrow/.style={thick, -{Stealth[scale=1.2]}},
      doublearrow/.style={thick, {Stealth[scale=1.2]}-{Stealth[scale=1.2]}},
      feedback/.style={thick, dashed, {Stealth[scale=1.2]}-{Stealth[scale=1.2]}, color=gray!80}
    ]

    % Definindo os Nós (Caixas)
    \node (fase1) [stepbox] {\textbf{Fase 1}\\Delineamento e seleção da abordagem};
    \node (fase2) [stepbox, right=of fase1] {\textbf{Fase 2}\\Procedimentos de elicitação e documentação};
    \node (fase3) [stepbox, below=of fase2] {\textbf{Fase 3}\\Análise de rastreabilidade e cobertura};
    \node (fase4) [stepbox, left=of fase3] {\textbf{Fase 4}\\Levantamento da percepção da equipe};

    % Desenhando as setas principais
    \draw [arrow] (fase1) -- (fase2);
    \draw [doublearrow] (fase2) -- (fase3);
    \draw [arrow] (fase3) -- (fase4);

    \end{tikzpicture}

    \vspace{0.5em}
    {\footnotesize \textbf{Fonte:} Elaborado pelo autor.}
  \end{figure}
\end{frame}
